\documentclass[10pt,a4paper]{amsart}
\usepackage[margin=19mm]{geometry}
\usepackage{amsmath,amssymb,mathtools}
\usepackage[scr=rsfs,cal=euler]{mathalfa}
\usepackage{xfrac}
\usepackage[unicode,hidelinks,bookmarks=false]{hyperref}
\newcommand{\ie}{i.e.\ }
\newcommand{\cf}{cf.\ }
\newcommand{\resp}{respectively\ }
\newcommand{\Subsection}{\subsection}
\providecommand{\nobreakdash}{\nobreak\hbox{-}}
\setlength{\emergencystretch}{2em}
\multlinegap0pt
\usepackage{amssymb}
\usepackage{mathtools}
\usepackage{xcolor}
\usepackage{tikz}
\newtheorem{theorem}{Theorem}
\title{The longest-edge bisection algorithm may produce degenerating tetrahedra}
\author{Sergey Korotov}
\date{Source: arXiv:2608.23139v1 (CC BY 4.0)}
\begin{document}
\maketitle

\noindent\emph{Source and licence.} The longest-edge bisection algorithm may produce degenerating tetrahedra. Version arXiv:2608.23139v1 (CC BY 4.0), distributed by arXiv under the Creative Commons Attribution 4.0 International licence, http://creativecommons.org/licenses/by/4.0/, as recorded in the arXiv licence metadata for this exact version and shown on the abstract page https://arxiv.org/abs/2608.23139v1. Full licence notice: \texttt{licenses/arxiv-2608.23139-CC-BY-4.0.txt}.\par\smallskip

\noindent\textit{Editorial scope.} Counted unit: the article's theorem thm:failure (source0.tex lines 221--235). The article's complete original proof of it is delivered with it (source0.tex lines 237--262, one \texttt{proof} environment of 26 lines). No mathematics is written, repaired or re-derived: the statement and the proof are the article's own lines, in the article's own notation, and no retained line is substituted or re-flowed.  Retained uncounted as the source-local context the proof needs: lines 204--207; lines 215--219.  Nothing was omitted from any retained span, so the package carries no omission record.  No retained line contains a citation, so the wrapper carries no bibliography and no bibliography entry.  No retained line contains a figure environment, an image-inclusion command or drawing code, so no image is delivered and the package declares no asset dependency.  Editorial formatting only, in addition: an AMS \texttt{amsart} preamble that restates the article's own declarations and macros used by the retained lines, namely the environments \texttt{theorem} on separate counters, together with the standard packages those lines need.  The retained environments print in this document as Theorem 1 in order of appearance, whereas the article prints the same items with the numbering of its own full document.  The complete native source of the article as distributed in the arXiv e-print for this exact version is bundled unchanged as \texttt{sources/208394/source0.tex}.
\begin{equation}\label{eq:chain}
E(1)\longrightarrow O(1)\longrightarrow E(1/2)\longrightarrow O(1/2)
\longrightarrow E(1/4)\longrightarrow\cdots.
\end{equation}

\begin{equation}\label{eq:volume}
|E(a)|=\frac{\sqrt7}{12}a^2,
\qquad
h_{E(a)}^2=|AD|^2=a(2+a).
\end{equation}

\begin{theorem}[Three simultaneous failures]\label{thm:failure}
The admissible sequence \eqref{eq:chain} violates shape regularity, the tetrahedral minimum-angle condition, and the tetrahedral maximum-angle condition.  More precisely,
\begin{align}
\frac{|E_k|}{h_{E_k}^3}
&=\frac{\sqrt7}{12}\frac{\sqrt{a_k}}{(2+a_k)^{3/2}}
\longrightarrow0, \label{eq:regfail}\\
\delta(a_k)
&=\arctan\sqrt{\frac{8a_k}{7}}
\longrightarrow0, \label{eq:minfail}\\
\cos\Theta(a_k)
&=\frac{2a_k-7}{\sqrt{(4a_k+7)(8a_k+7)}}
\longrightarrow-1, \label{eq:maxfail}
\end{align}
where $\delta(a)$ is the interior dihedral angle of $E(a)$ at $CD$, and $\Theta(a)$ is the interior dihedral angle of $O(a)$ at $MC$.  Thus $\Theta(a_k)\to\pi$.
\end{theorem}

\begin{proof}
Formula \eqref{eq:regfail} follows directly from \eqref{eq:volume}; its right-hand side is asymptotic to
\[
\frac{\sqrt7}{12\,2^{3/2}}\,2^{-k/2}.
\]
Therefore no uniform positive lower bound for $|T|/h_T^3$ is possible.

To obtain \eqref{eq:minfail}, intersect $E(a)$ with a plane perpendicular to $CD$.  Since
\[
d(B,CD)=\frac{2|BCD|}{|CD|}=\sqrt{\frac{7a}{8}}
\]
and $A$ lies vertically a distance $a$ above $B$, the relevant cross-sectional angle is
\[
\delta(a)=\arctan\frac{a}{\sqrt{7a/8}}
=\arctan\sqrt{\frac{8a}{7}}.
\]
Hence an interior dihedral angle tends to zero.

For the angle at $MC$ in $O(a)$, let $u=C-M$ and project $B-M$ and $D-M$ onto the plane perpendicular to $u$.  Calling the projected vectors $p_B$ and $p_D$, direct calculation gives
\[
|p_B|^2=\frac{a(4a+7)}{4(a+4)},\qquad
|p_D|^2=\frac{a(8a+7)}{4(a+4)},\qquad
p_B\!\cdot p_D=\frac{a(2a-7)}{4(a+4)}.
\]
Their angle is the interior dihedral angle $\Theta(a)$, which proves \eqref{eq:maxfail}.  As $a\to0$, its cosine tends to $-1$, and therefore $\Theta(a)\to\pi$.
\end{proof}

\end{document}
